\chapter{\texorpdfstring{附\quad 录}{附录}} %这里多出来的代码是用来消除hyperref包引入的警告的，\texorpdfstring{}{}命令的作用是告诉hyperref包在生成PDF书签时使用第二个参数，而不是第一个参数中的特殊字符。

\begin{appendixenv}

\section{样本描述性统计}

样本基本信息的描述性统计结果如表\ref{tab:sample_demo}所示。

\begin{table}[htbp]
    \centering
    \caption{样本的描述性统计}\zihao{5}
    \begin{tabular}{llcc}
        \toprule
        \songtibf{变量} & \songtibf{分类} & \songtibf{频数} & \songtibf{频率\nolinebreak[3]（\%）} \\
        \midrule
        \multirow{2}{*}{性别} & 男 & 76 & 49.35 \\
        & 女 & 78 & 50.65 \\
        \midrule
        \multirow{5}{*}{年级} & 大一 & 41 & 26.62 \\
        & 大二 & 39 & 25.32 \\
        & 大三 & 21 & 13.64 \\
        & 大四 & 38 & 24.68 \\
        & 研究生及以上 & 15 & 9.74 \\
        \midrule
        \multirow{5}{*}{专业} & 管理学 & 13 & 8.44 \\
        & 理学 & 47 & 30.52 \\
        & 工学 & 64 & 41.56 \\
        & 文学 & 15 & 6.49 \\
        & 经济学 & 38 & 12.99 \\
        \midrule
        \multirow{2}{*}{是否有过创业经历} & 是 & 42 & 27.27 \\
        & 否 & 122 & 72.73 \\
        \midrule
        \multirow{2}{*}{是否有过社会志愿服务经历} & 是 & 132 & 85.71 \\
        & 否 & 22 & 14.29 \\
        \midrule
        \multirow{3}{*}{参与课程类型} & 理论型 & 68 & 44.16 \\
        & 实践型 & 53 & 34.42 \\
        & 融合型 & 33 & 21.43 \\
        \bottomrule
    \end{tabular}
    \label{tab:sample_demo}
\end{table}

\clearpage
\section{问卷调查}

\subsection{第一部分 基本信息}

\begin{enumerate}
    \item 您的性别是 \quad $\square$男 \quad $\square$女
    \item 您的年级是 \rule{4cm}{0.4pt}
    \item 您的专业类别是 \rule{4cm}{0.4pt} \\
    \nolinebreak[3]（1.管理学 \quad 2.理学 \quad 3.工学 \quad 4.文学 \quad 5.经济学 \quad 6.法学 \quad 7.医学 \quad 8.教育学 \quad 9.哲学 \quad 10.历史学 \quad 11.农学 \quad 12.艺术学）
    \item 您是否参与过创业项目或者有过创业经历？ \quad $\square$是 \quad $\square$否
    \item 您是否参与过志愿服务活动？ \quad $\square$是 \quad $\square$否
    \item 在您选修过的创新创业课程中，哪一类是您学习时间最长、参与程度最深的？\quad $\square$理论型 \quad $\square$实践型 \quad $\square$融合型
\end{enumerate}

\subsection{第二部分 课程感知量表}

\noindent 注：以下题项均采用李克特五级量表，1\nolinebreak[3]（非常不符合）→5\nolinebreak[3]（非常符合）。

\tablemacro{H}
            {理论型课程感知量表}
            {p{10cm}c}
            {\songtibf{题目} & \songtibf{非常不符合 $\rightarrow$ 非常符合} \\}
            {课程中老师讲授的创业理论知识系统全面 & 1 \quad 2 \quad 3 \quad 4 \quad 5 \\
            课程中的案例教学对我理解创业很有启发 & 1 \quad 2 \quad 3 \quad 4 \quad 5 \\
            课程让我对创业的基本流程有了清晰认知 & 1 \quad 2 \quad 3 \quad 4 \quad 5 \\
            整体上，我对这门理论型课程感到满意 & 1 \quad 2 \quad 3 \quad 4 \quad 5 \\}
            {tab:scale_theory}

\tablemacro{H}
            {实践型课程感知量表}
            {p{10cm}c}
            {\songtibf{题目} & \songtibf{非常不符合 $\rightarrow$ 非常符合} \\}
            {课程提供了丰富的实践操作机会 & 1 \quad 2 \quad 3 \quad 4 \quad 5 \\
            我有机会将课堂所学应用到实际项目中去 & 1 \quad 2 \quad 3 \quad 4 \quad 5 \\
            课程中的实践环节提升了我的创业动手能力 & 1 \quad 2 \quad 3 \quad 4 \quad 5 \\
            整体上，我对这门实践型课程感到满意 & 1 \quad 2 \quad 3 \quad 4 \quad 5 \\}
            {tab:scale_practice}

\tablemacro{H}
            {融合型课程感知量表}
            {p{10cm}c}
            {\songtibf{题目} & \songtibf{非常不符合 $\rightarrow$ 非常符合} \\}
            {课程将理论知识与实践操作很好地结合在一起 & 1 \quad 2 \quad 3 \quad 4 \quad 5 \\
            通过课程学习，我能够运用理论指导实践 & 1 \quad 2 \quad 3 \quad 4 \quad 5 \\
            理论与实践的结合让我对创业有了更全面的理解 & 1 \quad 2 \quad 3 \quad 4 \quad 5 \\
            整体上，我对这门融合型课程的质量感到满意 & 1 \quad 2 \quad 3 \quad 4 \quad 5 \\}
            {tab:scale_fusion}

\subsection{第三部分 自主动机量表}

\tablemacro{H}
            {自主动机量表}
            {p{10cm}c}
            {\songtibf{题目} & \songtibf{非常不符合 $\rightarrow$ 非常符合} \\}
            {我对创业本身很感兴趣 & 1 \quad 2 \quad 3 \quad 4 \quad 5 \\
            参与创业课程让我感受到愉快和充实 & 1 \quad 2 \quad 3 \quad 4 \quad 5 \\
            我想通过创业课程挑战自己、实现自我价值 & 1 \quad 2 \quad 3 \quad 4 \quad 5 \\
            我认为创业对未来的职业发展很重要 & 1 \quad 2 \quad 3 \quad 4 \quad 5 \\
            学习创业对我的未来发展很有帮助 & 1 \quad 2 \quad 3 \quad 4 \quad 5 \\}
            {tab:scale_autonomous}

\subsection{第四部分 外部动机量表}

\tablemacro{H}
            {外部动机量表}
            {p{10cm}c}
            {\songtibf{题目} & \songtibf{非常不符合 $\rightarrow$ 非常符合} \\}
            {我参与创业课程是因为不参与我会感到不安 & 1 \quad 2 \quad 3 \quad 4 \quad 5 \\
            我选修这门课是为了证明我自己的能力 & 1 \quad 2 \quad 3 \quad 4 \quad 5 \\
            我选修这门课主要是为了获得学分 & 1 \quad 2 \quad 3 \quad 4 \quad 5 \\
            我参与创业活动是为了获得奖励或证书 & 1 \quad 2 \quad 3 \quad 4 \quad 5 \\
            我不太清楚自己为什么会选修这门课 & 1 \quad 2 \quad 3 \quad 4 \quad 5 \\}
            {tab:scale_external}

\subsection{第五部分 创业意愿量表}

\tablemacro{H}
            {创业意愿量表}
            {p{10cm}c}
            {\songtibf{题目} & \songtibf{非常不符合 $\rightarrow$ 非常符合} \\}
            {我有强烈的意愿在未来创办自己的企业 & 1 \quad 2 \quad 3 \quad 4 \quad 5 \\
            我的职业目标之一是成为一名创业者 & 1 \quad 2 \quad 3 \quad 4 \quad 5 \\
            我会认真考虑从事创业相关的工作 & 1 \quad 2 \quad 3 \quad 4 \quad 5 \\
            我愿意用尽一切努力来创办和经营自己的企业 & 1 \quad 2 \quad 3 \quad 4 \quad 5 \\
            我正在为创业做各种准备\nolinebreak[3]（学习、实习、积累资源） & 1 \quad 2 \quad 3 \quad 4 \quad 5 \\
            即使存在风险，我也愿意尝试创业 & 1 \quad 2 \quad 3 \quad 4 \quad 5 \\
            如果机会合适，我会立刻开始创业 & 1 \quad 2 \quad 3 \quad 4 \quad 5 \\}
            {tab:scale_intention}

\end{appendixenv}

\clearpage
